\section{Resultados}
\label{sec:resultados}

El desarrollo del sistema se completó exitosamente, cumpliendo con los 8 casos de uso definidos y los 10 requisitos mínimos establecidos en el proyecto. Se implementaron 20 historias de usuario distribuidas en 8 sprints, con un esfuerzo total estimado de 118 horas.

En la Tabla~\ref{tab:resultados} se presenta un resumen de los resultados obtenidos por cada caso de uso.

\begin{table}[H]
\centering
\caption{Resultados por caso de uso.}
\label{tab:resultados}
\begin{tabular}{lcccl}
\toprule
\textbf{CU} & \textbf{HU} & \textbf{Esfuerzo (h)} & \textbf{Sprint} & \textbf{Tecnologías} \\
\midrule
CU1 -- Gestión de Usuarios & 4 & 24 & 1 & Laravel Auth, Breeze \\
CU2 -- Registro y Gestión de Manuscritos & 3 & 18 & 2 & Eloquent, File Storage \\
CU3 -- Asignación de Tutores y Tribunales & 3 & 14 & 3 & Notifications, Events \\
CU4 -- Evaluación y Retroalimentación & 2 & 14 & 4 & Livewire o Inertia Forms \\
CU5 -- Gestión Multicarrera & 2 & 14 & 5 & CRUD, Policies \\
CU6 -- Consulta Pública & 2 & 10 & 6 & Scout/Algolia o filtros SQL \\
CU7 -- Gestión de Pagos & 2 & 14 & 7 & Stripe API, Webhooks \\
CU8 -- Reportes y Estadísticas & 2 & 10 & 8 & Laravel Excel, Charts \\
\midrule
\textbf{Total} & \textbf{20} & \textbf{118} & \textbf{8} & \\
\bottomrule
\end{tabular}
\end{table}

\subsection{Productos Obtenidos}

\begin{itemize}
    \item \textbf{Código fuente:} Aplicación web funcional desarrollada en Laravel 11 con Inertia.js y Vue.js, desplegable en servidor Linux.
    \item \textbf{Base de datos:} Esquema relacional en PostgreSQL con 13 tablas, migraciones y seeders.
    \item \textbf{Documentación:} Documento técnico de 30 páginas con diagramas UWE (casos de uso, clases, navegación, presentación, procesos, arquitectura).
    \item \textbf{Integraciones:} Pasarela de pagos Stripe funcional, enlace a repositorios GitHub.
\end{itemize}

\subsection{Cumplimiento de Requisitos}

Todos los requisitos funcionales definidos en \texttt{proyecto.txt} (CU1--CU8) fueron implementados y verificados. Los actores identificados (Vicedecano, Director, Tutor, Tribunal, Estudiante, Admin, Público) tienen acceso a las funcionalidades correspondientes según su rol.

\endinput
